\section{Brechas y oportunidades de investigación}

La primera brecha es la \textbf{escasez de datos clínicamente anotados}. Los mejores resultados cuantitativos suelen provenir de grandes corpus con etiquetas simplificadas, mientras que los recursos más útiles para interpretación clínica, como los alineados con C-SSRS o DSM-5, son considerablemente más pequeños y costosos de construir \citep{gaur2019severity, bao2025redsm5}. Esto limita la posibilidad de entrenar modelos potentes sin sacrificar granularidad clínica.

La segunda brecha es la \textbf{desconexión entre desempeño y explicabilidad}. Muchos artículos reportan exactitud, F1 o AUC competitivos, pero ofrecen explicaciones débiles sobre qué evidencia sustenta la decisión del modelo. Desde la dimensión de efectividad predictiva, esto puede parecer suficiente; sin embargo, desde la dimensión de interpretabilidad clínica, la limitación es central, porque una predicción útil debería poder relacionarse con síntomas, racionales textuales o indicadores observables para disminuir el riesgo de uso acrítico del sistema \citep{montejo2024survey, mirtaheri2024tcn, shuvo2024bigru}.

La tercera brecha es la \textbf{limitada validación ecológica}. Varios estudios trabajan con Reddit o Twitter como fuente principal y, aunque ello es metodológicamente valioso, todavía existe poca evidencia de generalización hacia otros idiomas, otras plataformas o contextos clínicos reales. Los artículos de revisión coinciden en señalar un sesgo fuerte hacia el inglés y hacia configuraciones experimentales controladas \citep{abdulsalam2024suicidal, montejo2024survey}.

La cuarta brecha es la \textbf{insuficiente integración temporal y multimodal}. El cambio de estado, la progresión de riesgo y la historia previa del usuario son clínicamente relevantes, pero no siempre quedan modelados cuando la tarea se formula como clasificación estática de una sola publicación. Del mismo modo, el campo ha explorado más texto que combinaciones con metadatos, imágenes o patrones de interacción. Varias revisiones identifican precisamente ese camino como una dirección prioritaria \citep{dechoudhury2016shifts, chatterjee2022multimodal, montejo2024survey}.

La quinta brecha es la \textbf{falta de criterios homogéneos para medir interpretabilidad clínica}. Algunos trabajos la asocian con categorías DSM-5, otros con niveles de severidad, otros con recuperación de evidencia y otros con análisis \textit{post hoc} mediante LLMs. Esta variedad es valiosa, pero también dificulta comparar enfoques y seleccionar un marco consistente para futuras fases del proyecto. En consecuencia, una oportunidad clara consiste en diseñar sistemas que combinen potencia predictiva con salidas explícitas y verificables, por ejemplo, síntomas, evidencia textual y nivel de riesgo. Este punto aporta una guía concreta para el reto, porque sugiere que la comparación entre modelos debe considerar no solo quién clasifica mejor, sino también quién explica mejor.
